% ==================== 文档类 ====================
\documentclass[12pt, a4paper]{article}
% ==================== 中文支持 ====================
\usepackage[UTF8]{ctex} 
% ==================== 页面设置 ====================
\usepackage{geometry}
\geometry{
    left=2.5cm,
    right=2.5cm,
    top=2.5cm,
    bottom=2.5cm
}

% ==================== 数学公式 ====================
\usepackage{amsmath} % 基础数学环境 
\usepackage{amssymb} % 数学符号
\usepackage{amsthm}  % 定理环境

% ==================== 图表支持 ====================
\usepackage{graphicx}  % 插入图片
\usepackage{float}     % 浮动体位置控制
\usepackage{subfigure} % 子图支持
\usepackage{booktabs}  % 三线表
\usepackage{multirow}  % 表格合并行
\usepackage{array}     % 表格增强

% ==================== 代码高亮 ====================
\usepackage{listings}  % 代码块
\usepackage{xcolor}    % 颜色支持

% 代码样式设置
\lstset{
    basicstyle=\ttfamily\small,
    keywordstyle=\color{blue},
    commentstyle=\color{gray},
    stringstyle=\color{red},
    numbers=left,
    numberstyle=\tiny\color{gray},
    stepnumber=1,
    numbersep=8pt,
    showstringspaces=false,
    breaklines=true,
    frame=single,
    backgroundcolor=\color{gray!10}
}

% ==================== 超链接 ====================
\usepackage{hyperref}
\hypersetup{
    colorlinks=true,
    linkcolor=blue,
    citecolor=blue,
    urlcolor=blue
}

% ==================== 参考文献 ====================
\usepackage{cite}

% ==================== 实用宏包 ====================
\usepackage{enumerate} % 自定义列表
\usepackage{enumitem}  % 列表增强
\usepackage{caption}   % 图表标题
\usepackage{titlesec}  % 标题格式

% ==================== 文档信息 ====================
\title{文档标题}
\author{作者姓名
\date{\today}

% ==================== 正文开始 ====================
\begin{document}
\maketitle
\section{皮尔逊相关系数 (Pearson Correlation Coefficient)}
皮尔逊相关系数是一种衡量两个变量之间线性相关程度的统计指标。其定义为两个变量的协方差与其标准差乘积的比值，通常用字母 $r$ 表示。
\subsection{数学定义}
对于样本数据集 $(x_1, y_1), (x_2, y_2), \dots, (x_n, y_n)$，样本皮尔逊相关系数 $r$ 的计算公式如下：
\begin{equation}
r = \frac{\sum_{i=1}^{n} (x_i - \bar{x})(y_i - \bar{y})}{\sqrt{\sum_{i=1}^{n} (x_i - \bar{x})^2 \sum_{i=1}^{n} (y_i - \bar{y})^2}}
\end{equation}

其中：
\begin{itemize}
    \item $\bar{x}$ 和 $\bar{y}$ 分别是变量 $x$ 和 $y$ 的样本均值。
    \item 分子表示变量 $x$ 和 $y$ 的协方差。
    \item 分母是变量 $x$ 和 $y$ 标准差的乘积。
\end{itemize}

\subsection{取值意义}

皮尔逊相关系数的取值范围为 $[-1, 1]$：

\begin{enumerate}
    \item \textbf{$r > 0$}：表示两个变量正相关。
    \item \textbf{$r < 0$}：表示两个变量负相关。
    \item \textbf{$|r| = 1$}：表示变量之间存在完全线性的关系。
    \item \textbf{$r = 0$}：表示变量之间不存在线性相关关系（但可能存在非线性关系）。
\end{enumerate}

\end{document}